
\documentclass{article}
\usepackage{colortbl}
\usepackage{makecell}
\usepackage{multirow}
\usepackage{supertabular}

\begin{document}

\newcounter{utterance}

\centering \large Interaction Transcript for game `cladder', experiment `full\_v1.5\_default', episode 1876 with selfrefine{-}Qwen3.5{-}27B.
\vspace{24pt}

{ \footnotesize  \setcounter{utterance}{1}
\setlength{\tabcolsep}{0pt}
\begin{supertabular}{c@{$\;$}|p{.15\linewidth}@{}p{.15\linewidth}p{.15\linewidth}p{.15\linewidth}p{.15\linewidth}p{.15\linewidth}}
    \# & \multicolumn{2}{c}{Player} && \multicolumn{2}{c}{Game Master} \\
    \hline

    \theutterance \stepcounter{utterance}  
    & & & \multicolumn{4}{p{0.6\linewidth}}{
        \cellcolor[rgb]{0.9,0.9,0.9}{
            \makecell[{{p{\linewidth}}}]{
                \texttt{\tiny{[P1$\langle$GM]}}
                \texttt{You are an expert in causal inference. The following question is not a typical commonsense query, but rather a meticulously designed question created by a professor specializing in causal inference, intended to assess the students' mastery of the course content.} \\
\texttt{Imagine a self{-}contained, hypothetical world with only the following conditions, and without any unmentioned factors or causal relationships: Medication has a direct effect on blood pressure and heart condition. Blood pressure has a direct effect on heart condition. For people not taking any medication, the probability of healthy heart is 64\%. For people taking medication, the probability of healthy heart is 21\%. For people taking medication, would it be more likely to see healthy heart if the person had taken no medication?} \\
\texttt{Start your answer with 'Yes' or 'No', followed by additional reasoning or evidence to support your explanation.} \\
            }
        }
    }
     \\ \\

    \theutterance \stepcounter{utterance}  
    & \multicolumn{4}{p{0.6\linewidth}}{
        \cellcolor[rgb]{0.561,0.737,0.561}{
            \makecell[{{p{\linewidth}}}]{
                \texttt{\tiny{[P1$\rangle$GM]}}
                \texttt{No} \\
\\ 
\texttt{**Reasoning:**} \\
\\ 
\texttt{1.  **Analyze the Causal Structure:**} \\
\texttt{    The prompt explicitly defines a self{-}contained world with specific causal links:} \\
\texttt{    *   Medication \$\textbackslash{}to\$ Blood Pressure} \\
\texttt{    *   Medication \$\textbackslash{}to\$ Heart Condition} \\
\texttt{    *   Blood Pressure \$\textbackslash{}to\$ Heart Condition} \\
\texttt{    Crucially, the prompt states there are **no unmentioned factors or causal relationships**. This explicitly rules out the real{-}world confounding scenario where "sicker patients are more likely to be prescribed medication" (i.e., there is no hidden variable \$U\$ causing both Medication and Heart Condition).} \\
\\ 
\texttt{2.  **Interpret the Observed Probabilities:**} \\
\texttt{    *   \$P(\textbackslash{}text\{Healthy Heart\} \textbackslash{}mid \textbackslash{}text\{No Medication\}) = 64\textbackslash{}\%\$} \\
\texttt{    *   \$P(\textbackslash{}text\{Healthy Heart\} \textbackslash{}mid \textbackslash{}text\{Medication\}) = 21\textbackslash{}\%\$} \\
\texttt{    In a world without confounding, the observed association directly reflects the causal effect. The data shows that the group taking medication has a significantly lower probability of a healthy heart compared to the group not taking it.} \\
\\ 
\texttt{3.  **Evaluate the Counterfactual:**} \\
\texttt{    The question asks: "For people taking medication, would it be more likely to see healthy heart if the person had taken no medication?"} \\
\texttt{    *   This asks for the counterfactual outcome \$Y\_\{\textbackslash{}text\{no med\}\}\$ for the treated group, compared to their actual outcome \$Y\_\{\textbackslash{}text\{med\}\}\$.} \\
\texttt{    *   Since there is no confounding, the causal effect of medication is negative (it reduces the probability of a healthy heart from 64\% to 21\%).} \\
\texttt{    *   Therefore, if a person currently taking medication were to stop taking it (the counterfactual "no medication" scenario), their probability of having a healthy heart would increase (moving from the 21\% baseline of the treated group toward the 64\% baseline of the untreated group).} \\
\texttt{    *   Wait, let's re{-}read the specific phrasing carefully: "would it be more likely to see healthy heart if the person had taken no medication?"} \\
\texttt{    *   Current state (Medication): 21\% chance of healthy heart.} \\
\texttt{    *   Counterfactual state (No Medication): Since the causal effect of medication is harmful (lowering the rate from 64\% to 21\%), removing the medication would restore the probability to the higher level (64\%).} \\
\texttt{    *   Therefore, the probability of a healthy heart *without* medication (64\%) is indeed **higher** than *with* medication (21\%).} \\
\\ 
\texttt{    *Correction on the previous logic step:* The previous "No" draft was based on the assumption that the medication was beneficial. But the data says \$P(H|M) = 21\textbackslash{}\%\$} \\
            }
        }
    }
    & & \\ \\

    \theutterance \stepcounter{utterance}  
    & & & \multicolumn{2}{p{0.3\linewidth}}{
        \cellcolor[rgb]{0.9,0.9,0.9}{
            \makecell[{{p{\linewidth}}}]{
                \texttt{\tiny{[GM$|$GM]}}
                \texttt{No} \\
            }
        }
    }
    & & \\ \\

    \theutterance \stepcounter{utterance}  
    & & & \multicolumn{2}{p{0.3\linewidth}}{
        \cellcolor[rgb]{0.9,0.9,0.9}{
            \makecell[{{p{\linewidth}}}]{
                \texttt{\tiny{[GM$|$GM]}}
                \texttt{yes} \\
            }
        }
    }
    & & \\ \\

    \theutterance \stepcounter{utterance}  
    & & & \multicolumn{2}{p{0.3\linewidth}}{
        \cellcolor[rgb]{0.9,0.9,0.9}{
            \makecell[{{p{\linewidth}}}]{
                \texttt{\tiny{[GM$|$GM]}}
                \texttt{game\_result = LOSE} \\
            }
        }
    }
    & & \\ \\

\end{supertabular}
}

\end{document}
